% =========================================================
% Section: One-Sided Derivatives
% =========================================================

\begin{tcolorbox}[colback=gray!8, colframe=gray!40, title=\textbf{One-Sided Derivatives --- toolkit}]
\textbf{Concept family:} Left-hand and right-hand derivatives, and their relation to full differentiability.

\medskip
\textbf{Atomic items in this section:}
\begin{itemize}
  \item Definition: Left-Hand Derivative
  \item Definition: Right-Hand Derivative
  \item Theorem: Differentiability and One-Sided Derivatives
\end{itemize}
\end{tcolorbox}

% ---------------------------------------------------------

\begin{tcolorbox}[colback=defbox, colframe=defborder]
\begin{definition}[Left-Hand Derivative]
\label{def:left-hand-derivative}
Let $I \subseteq \mathbb{R}$ be an interval, let $f : I \to \mathbb{R}$, and let $c \in I$ be such that $c$ is a left-hand limit point of $I$. We say that $f$ has a \emph{left-hand derivative} at $c$ if the limit
\[
\lim_{x \to c^-} \frac{f(x) - f(c)}{x - c}
\]
exists. In that case, the value of this limit is denoted $f'_-(c)$:
\[
f'_-(c) := \lim_{x \to c^-} \frac{f(x) - f(c)}{x - c}.
\]
\end{definition}
\end{tcolorbox}

\begin{remark*}[Standard quantified statement]
\[
\exists L \in \mathbb{R}\;\forall \varepsilon > 0\;\exists \delta > 0\;\forall x \in I
\Bigl(-({\delta}) < x - c < 0 \implies \left|\frac{f(x)-f(c)}{x-c} - L\right| < \varepsilon\Bigr).
\]
In this case $f'_-(c) = L$.
\end{remark*}

\begin{remark*}[Definition predicate reading]
\[
\operatorname{LeftDerivativeAt}(f, c, L)
\;\Longleftrightarrow\;
\operatorname{LeftLimit}\!\left(\lambda x.\,\frac{f(x)-f(c)}{x-c},\, c,\, L\right).
\]
\end{remark*}

\begin{remark*}[Negated quantified statement]
\[
\forall L \in \mathbb{R}\;\exists \varepsilon > 0\;\forall \delta > 0\;\exists x \in I
\Bigl(-\delta < x - c < 0 \;\land\; \left|\frac{f(x)-f(c)}{x-c} - L\right| \ge \varepsilon\Bigr).
\]
\end{remark*}

\begin{remark*}[Negation predicate reading]
\[
\forall L \in \mathbb{R}\;\lnot\operatorname{LeftDerivativeAt}(f, c, L).
\]
No finite $L$ is the left-hand limit of the difference quotient.
\end{remark*}

\begin{remark*}[Interpretation]
The left-hand derivative $f'_-(c)$ is the limit of the difference quotient as $x$ approaches $c$ \emph{strictly from the left}. It measures the instantaneous rate of change of $f$ at $c$ as seen from the left side. The right-hand derivative $f'_+(c)$ is the symmetric object from the right. At interior points of an interval, both one-sided derivatives are defined, and the full derivative $f'(c)$ exists if and only if they agree. At endpoints, only the inward one-sided derivative is accessible, and it serves as the boundary derivative. The canonical example: for $f(x) = |x|$ at $c = 0$, $f'_-(0) = -1$ and $f'_+(0) = +1$; since these disagree, $f'(0)$ does not exist.
\end{remark*}

\begin{remark*}[Dependencies]
\begin{itemize}
  \item \hyperref[def:derivative-at-a-point]{Definition: Derivative at a Point}
\end{itemize}
\end{remark*}

% ---------------------------------------------------------

\begin{tcolorbox}[colback=defbox, colframe=defborder]
\begin{definition}[Right-Hand Derivative]
\label{def:right-hand-derivative}
Let $I \subseteq \mathbb{R}$ be an interval, let $f : I \to \mathbb{R}$, and let $c \in I$ be such that $c$ is a right-hand limit point of $I$. We say that $f$ has a \emph{right-hand derivative} at $c$ if the limit
\[
\lim_{x \to c^+} \frac{f(x) - f(c)}{x - c}
\]
exists. In that case, the value of this limit is denoted $f'_+(c)$:
\[
f'_+(c) := \lim_{x \to c^+} \frac{f(x) - f(c)}{x - c}.
\]
\end{definition}
\end{tcolorbox}

\begin{remark*}[Standard quantified statement]
\[
\exists L \in \mathbb{R}\;\forall \varepsilon > 0\;\exists \delta > 0\;\forall x \in I
\Bigl(0 < x - c < \delta \implies \left|\frac{f(x)-f(c)}{x-c} - L\right| < \varepsilon\Bigr).
\]
In this case $f'_+(c) = L$.
\end{remark*}

\begin{remark*}[Definition predicate reading]
\[
\operatorname{RightDerivativeAt}(f, c, L)
\;\Longleftrightarrow\;
\operatorname{RightLimit}\!\left(\lambda x.\,\frac{f(x)-f(c)}{x-c},\, c,\, L\right).
\]
\end{remark*}

\begin{remark*}[Negated quantified statement]
\[
\forall L \in \mathbb{R}\;\exists \varepsilon > 0\;\forall \delta > 0\;\exists x \in I
\Bigl(0 < x - c < \delta \;\land\; \left|\frac{f(x)-f(c)}{x-c} - L\right| \ge \varepsilon\Bigr).
\]
\end{remark*}

\begin{remark*}[Negation predicate reading]
\[
\forall L \in \mathbb{R}\;\lnot\operatorname{RightDerivativeAt}(f, c, L).
\]
\end{remark*}

\begin{remark*}[Dependencies]
\begin{itemize}
  \item \hyperref[def:derivative-at-a-point]{Definition: Derivative at a Point}
  \item \hyperref[def:left-hand-derivative]{Definition: Left-Hand Derivative}
\end{itemize}
\end{remark*}

% ---------------------------------------------------------

\begin{theorem}[Differentiability and One-Sided Derivatives]
\label{thm:differentiability-and-one-sided-derivatives}
Let $I \subseteq \mathbb{R}$ be an open interval and let $f : I \to \mathbb{R}$. Then $f$ is differentiable at $c \in I$ if and only if both one-sided derivatives $f'_-(c)$ and $f'_+(c)$ exist and are equal. In that case,
\[
f'(c) = f'_-(c) = f'_+(c).
\]
\hyperref[prf:differentiability-and-one-sided-derivatives]{Go to proof.}
\end{theorem}

\begin{remark*}[Standard quantified statement]
\[
\operatorname{DifferentiableAt}(f, c)
\;\iff\;
\Bigl(\operatorname{LeftDerivativeAt}(f,c,L) \;\land\; \operatorname{RightDerivativeAt}(f,c,L)\Bigr)
\;\text{for some } L \in \mathbb{R}.
\]
\end{remark*}

\begin{remark*}[Negated quantified statement]
\[
\lnot\operatorname{DifferentiableAt}(f,c)
\;\iff\;
\lnot\Bigl(\operatorname{LeftDerivativeAt}(f,c,L) \;\land\; \operatorname{RightDerivativeAt}(f,c,L)\Bigr)
\;\text{for all } L \in \mathbb{R}.
\]
That is: either a one-sided derivative fails to exist, or both exist but are unequal.
\end{remark*}

\begin{remark*}[Interpretation]
This theorem decomposes differentiability at an interior point into two one-sided conditions. It is the standard tool for proving non-differentiability: to show $f$ is not differentiable at $c$, compute $f'_-(c)$ and $f'_+(c)$ and observe they disagree. For $f(x) = |x|$ at $c = 0$: $f'_-(0) = -1 \neq +1 = f'_+(0)$, so $f$ is not differentiable at $0$. The theorem also provides the correct definition of differentiability on a closed interval $[a,b]$: differentiable on the open interior $(a,b)$ in the full sense, plus $f'_+(a)$ and $f'_-(b)$ both exist at the endpoints.
\end{remark*}

\begin{remark*}[Dependencies]
\begin{itemize}
  \item \hyperref[def:left-hand-derivative]{Definition: Left-Hand Derivative}
  \item \hyperref[def:right-hand-derivative]{Definition: Right-Hand Derivative}
  \item \hyperref[def:derivative-at-a-point]{Definition: Derivative at a Point}
\end{itemize}
\end{remark*}
